% Part: first-order-logic
% Chapter: syntax-and-semantics
% Section: covered-structures

\documentclass[../../../include/open-logic-section]{subfiles}

\begin{document}

\olfileid{fol}{syn}{cov}

\olsection{మొదటిస్థాయి విధేయ భాషల కోసం ఆవృత \printtoken{P}{structure}}

\begin{explain}
ఏ \tetoken{చరాలనూ}{!!{variable}s} కలిగి లేని పదాన్ని \emph{సంవృత} పదం అంటామని గుర్తుచేసుకోండి.
\end{explain}

\begin{defn}[సంవృత పదాల \tetoken{విలువ}{!!^{value}}]
$t$ అనేది భాష~$\Lang L$లోని సంవృత పదం, $\Struct M$ అనేది
$\Lang L$కు ఒక \tetoken{నిర్మాణం}{!!{structure}} అయితే, దాని \emph{\tetoken{విలువ}{!!{value}}}~$\Value{t}{M}$ను
కింది విధంగా నిర్వచిస్తాం:
\begin{enumerate}
\item $t$ కేవలం \tetoken{స్థిరసంకేతం}{!!{constant}}~$c$ అయితే, $\Value{c}{M} = \Assign{c}{M}$.
\item $t$ రూపం $\Atom{f}{t_1, \ldots, t_n}$ అయితే,
  \[
  \Value{t}{M} = \Assign{f}{M}(\Value{t_1}{M}, \ldots,
  \Value{t_n}{M}).
  \]
\end{enumerate}
\end{defn}

\begin{defn}[ఆవృత \tetoken{నిర్మాణం}{!!{structure}}]
ఒక \tetoken{నిర్మాణంలో}{!!{structure}} వ్యక్తి క్షేత్రంలోని ప్రతి మూలకం ఏదో ఒక సంవృత
పదం యొక్క \tetoken{విలువ}{!!{value}} అయితే, దానిని \emph{ఆవృతమైనది} అంటాం.
\end{defn}

\begin{ex}
~$\Lang L$ అనేది $\Obj{zero}$, $\Obj{one}$, $\Obj{two}$, \dots, అనే
\tetoken{స్థిరసంకేతాలు}{!!{constant}s}, ద్విస్థాన \tetoken{విధేయ సంకేతం}{!!{predicate}}~$<$, ద్విస్థాన \tetoken{ప్రమేయ సంకేతాలు}{!!{function}s}
$+$, $\times$ గల భాష అనుకుందాం. అప్పుడు~$\Lang L$కు \tetoken{ఒక నిర్మాణం}{!!a{structure}}~$\Struct
M$ ఇలా ఉంటుంది: దాని వ్యక్తి క్షేత్రం $\Domain M = \{0, 1, 2, \ldots
\}$; దాని నిర్దేశాలు $\Assign{\Obj{zero}}{M} = 0$,
$\Assign{\Obj{one}}{M} = 1$, $\Assign{\Obj{two}}{M} = 2$, ఇలా
కొనసాగుతాయి. ద్విస్థాన సంబంధ సంకేతం~$<$ కోసం, $\Assign{<}{M}$ అనేది
$c_1$, ~$c_2$ కంటే చిన్నదిగా ఉండే అన్ని జంటలు
$\tuple{c_1, c_2} \in \Domain{M}^2$ సమితి: ఉదాహరణకు,
$\tuple{1, 3} \in \Assign{<}{M}$, కానీ
$\tuple{2, 2} \notin \Assign{<}{M}$. ద్విస్థాన \tetoken{ప్రమేయ సంకేతం}{!!{function}}~$+$కు
$\Assign{+}{M}$ను సాధారణ విధంగా నిర్వచించండి---ఉదాహరణకు,
$\Assign{+}{M}(2,3)$ అనేది ~$5$కు పంపుతుంది; ద్విస్థాన
\tetoken{ప్రమేయ సంకేతం}{!!{function}}~$\times$కు కూడా అలాగే చేయండి. అందువల్ల $\Obj{four}$ యొక్క
\tetoken{విలువ}{!!{value}} కేవలం~$4$; $\times(\Obj{two},
+(\Obj{three},\Obj{zero}))$ యొక్క \tetoken{విలువ}{!!{value}} (లేదా మధ్యస్థ సంజ్ఞావిధానంలో,
$\Obj{two} \times (\Obj{three} + \Obj{zero})$) ఇది:\sourcecorrection{OLTEFOLSEM-002}{లెక్కలో వరుసగా వచ్చిన రెండు సమానతా చిహ్నాల్లో మొదటిదాన్ని తొలగించాం.}
\begin{multline*}
\Value{\times(\Obj{two}, +(\Obj{three},\Obj{zero}))}{M}\\
\begin{aligned}
& =\Assign{\times}{M}(\Value{\Obj{two}}{M}, \Value{+(\Obj{three}, \Obj{zero})}{M})\\
& = \Assign{\times}{M}(\Value{\Obj{two}}{M}, \Assign{+}{M}(\Value{\Obj{three}}{M},
\Value{\Obj{zero}}{M})) \\
& = \Assign{\times}{M}(\Assign{\Obj{two}}{M}, \Assign{+}{M}(\Assign{\Obj{three}}{M},
\Assign{\Obj{zero}}{M})) \\
& = \Assign{\times}{M}(2, \Assign{+}{M}(3, 0)) \\
& = \Assign{\times}{M}(2, 3) \\
& = 6
\end{aligned}
\end{multline*}
\end{ex}

\begin{prob}
అంకగణిత ప్రామాణిక నమూనా~$\Struct N$ ఆవృతమైనదా? వివరించండి.
\end{prob}

\end{document}
